% RadixGates 实验记录 · 展示幻灯片
% 用 XeLaTeX 编译:
%   latexmk -xelatex 实验记录-slides.tex

\documentclass[aspectratio=169,11pt]{beamer}

\usepackage[UTF8,fontset=macnew]{ctex}   % 中文支持（macOS 苹方/宋体）

\usetheme{Madrid}
\usecolortheme{default}
\setbeamertemplate{navigation symbols}{}
\setbeamertemplate{footline}[frame number]
\setbeamertemplate{itemize items}[circle]

\usepackage{booktabs}
\usepackage{tikz}
\usetikzlibrary{positioning,arrows.meta,fit,calc,backgrounds,shapes.geometric}
\usepackage{listings}
\usepackage{xcolor}

% ---- 代码高亮 ----
\definecolor{gokw}{RGB}{7,102,153}
\definecolor{gostr}{RGB}{170,55,55}
\definecolor{gocom}{RGB}{90,140,90}
\definecolor{gobg}{RGB}{247,247,247}
\definecolor{rgblue}{RGB}{20,90,150}
\definecolor{rggreen}{RGB}{35,130,80}
\definecolor{rgorange}{RGB}{200,110,20}
\definecolor{rgred}{RGB}{180,50,50}
\definecolor{rgpurple}{RGB}{120,60,150}

\lstdefinelanguage{Go}{
  morekeywords={break,default,func,interface,select,case,defer,go,map,struct,
    chan,else,goto,package,switch,const,fallthrough,if,range,type,continue,
    for,import,return,var,nil,true,false,make,new,append,len},
  sensitive=true,
  morecomment=[l]{//},
  morecomment=[s]{/*}{*/},
  morestring=[b]",
  morestring=[b]`,
}
\lstset{
  language=Go,
  backgroundcolor=\color{gobg},
  basicstyle=\ttfamily\scriptsize,
  keywordstyle=\color{gokw}\bfseries,
  stringstyle=\color{gostr},
  commentstyle=\color{gocom}\itshape,
  showstringspaces=false,
  columns=fullflexible,
  keepspaces=true,
  breaklines=true,
  frame=single,
  framesep=4pt,
  rulecolor=\color{gray!40},
  aboveskip=4pt, belowskip=2pt,
}
\newcommand{\gc}[1]{\texttt{\small #1}}   % 行内代码
\newcommand{\fpath}[1]{{\color{rgblue}\texttt{\small #1}}}
\newcommand{\cnum}[1]{\tikz[baseline=(c.base)]{\node[draw,circle,inner sep=1.2pt,line width=0.5pt](c){\scriptsize #1};}}

% ---- 通用 tikz 样式 ----
\tikzset{
  cli/.style={draw,rounded corners,fill=gray!10,align=center,font=\scriptsize,minimum height=7mm},
  gw/.style={draw,rounded corners,fill=rgblue!14,align=center,font=\scriptsize,minimum height=7mm},
  gpu/.style={draw,rounded corners,fill=rggreen!12,align=center,font=\scriptsize,minimum height=7mm},
  mod/.style={draw,rounded corners,fill=rgblue!10,align=center,font=\scriptsize,minimum height=8mm,minimum width=26mm},
  wstep/.style={draw,rounded corners,fill=rgblue!8,align=center,font=\scriptsize,minimum height=7mm,minimum width=17mm},
  dec/.style={draw,diamond,aspect=2,fill=rgorange!14,align=center,font=\scriptsize,inner sep=1pt},
  ar/.style={-{Latex[length=2mm]},thick},
}

\title[RadixGates 实验记录]{RadixGates 实验记录}
\subtitle{Vast.ai 4$\times$ RTX 4090 $\cdot$ 复现与验证}
\author{AI-Infra 课程}
\date{}

\begin{document}

% ==================================================================
\begin{frame}
  \titlepage
  \vspace{-1em}
  \small\centering
  在 Vast.ai 的 4$\times$ RTX 4090 实例上，跑通\textbf{流式 SSE / PD / DP / 企业压测 / TP / EP} 六个实验；\\
  真实记录\textbf{遇到的四个坑}与\textbf{与论文基准的对照}。
\end{frame}

% ------------------------------------------------------------------
\begin{frame}{一、机器环境（Vast.ai VM）}
\small
\begin{tabular}{@{}p{3.6cm}p{9.0cm}@{}}
\toprule
\textbf{项目} & \textbf{配置} \\
\midrule
GPU        & 4$\times$ NVIDIA RTX 4090（每卡 24GB，可用约 23GB） \\
\addlinespace
卡间互联   & \textbf{PHB（PCIe，不是 NVLink）} \\
\addlinespace
CPU / 内存 & 46 核 ｜ 197GB 内存 ｜ 126GB 系统盘（初始 119GB 空闲） \\
\addlinespace
容器       & Docker 28.1.1 ｜ docker compose v2.35.1 \\
\addlinespace
驱动 / CUDA & NVIDIA 驱动 575.51.03（\textbf{最高支持 CUDA 12.9}） \\
\addlinespace
其它       & git / python3 / curl 预装；端口 8080 空闲（无 Jupyter 抢占） \\
\bottomrule
\end{tabular}
\vfill
\centering\scriptsize
这是能跑 Docker 的 VM 实例（非受限 container template），所以走 Docker 路径部署。
\end{frame}

% ------------------------------------------------------------------
\begin{frame}{二、时间线概览：从连接到六实验跑通}
\small
\begin{enumerate}
  \item \textbf{第 0 步} 连接 + 环境体检：4$\times$4090 / 互联 PHB(PCIe) / Docker OK，GPU 直通可见。
  \item \textbf{第 1 步} 传代码 + 构建纯 Go 网关镜像（\textbf{95MB}）+ 后台预拉 SGLang 镜像。
  \item \textbf{起 DP} 撞坑 \cnum{1}（镜像 CUDA$\geq$13 vs 驱动 12.9）、\cnum{2}（\gc{--enable-prefix-caching} 废弃）$\to$ 修复。
  \item \textbf{第 4–7 步} \textcolor{rggreen}{\textbf{DP=4 $\to$ PD 分离 $\to$ 流式 SSE $\to$ 企业问答}} 依次跑通 ✅
  \item \textbf{起 TP} 撞坑 \cnum{3}（多卡 NCCL）$\to$ 加 \gc{ipc:host} + \gc{NCCL\_P2P\_DISABLE} $\to$ \textcolor{rggreen}{\textbf{TP=4 ✅}}
  \item \textbf{起 EP} 撞坑 \cnum{4}（\gc{--enable-expert-parallel} 废弃）$\to$ 改 \gc{--ep-size 4} $\to$ \textcolor{rggreen}{\textbf{EP=4 ✅}}
\end{enumerate}
\vfill
\centering\scriptsize 所有 compose 共享一个 \gc{hf\_cache} 卷；每个大模型实验前清卷腾盘（Qwen32B$\approx$64GB、MoE$\approx$28GB）。
\end{frame}

% ------------------------------------------------------------------
\begin{frame}[fragile]{第 0 步 + 第 1 步：连接与准备}
\textbf{第 0 步 — 连接 + 只读环境体检}：加载密钥后 SSH（\gc{-L} 把网关 8080 转发到本地）。
\begin{lstlisting}[language=bash]
ssh-add ~/.ssh/id_ed25519
ssh -p 22 root@203.0.113.10 -L 8080:localhost:8080
nvidia-smi ; docker version ; nvidia-smi topo -m   # 4x4090 / Docker 28.1.1 / PHB
docker run --rm --gpus all nvidia/cuda:12.4.0-base-ubuntu22.04 nvidia-smi -L
\end{lstlisting}
\textbf{第 1 步 — 传代码 + 构建网关镜像}（纯 Go、无 CGO，镜像仅 \textbf{95MB}）。
\begin{lstlisting}[language=bash]
docker pull lmsysorg/sglang:latest            # 后台预拉，边下边干
cd /root/RadixGates && docker build -t unicoregpu2020/radixgates:latest .
pip3 install requests
\end{lstlisting}
\centering\scriptsize
关键点：所有 compose 共享一个 \gc{hf\_cache} 卷，同一模型只下载一次、跨实验复用。
\end{frame}

% ------------------------------------------------------------------
\begin{frame}[fragile]{问题 \cnum{1}：镜像 CUDA$\geq$13 vs 驱动 12.9}
\textbf{现象}：首次启动 DP=4，容器起不来，报错：
\begin{lstlisting}[language=]
nvidia-container-cli: requirement error: unsatisfied condition:
cuda>=13.0, please update your driver ... or use an earlier cuda container
\end{lstlisting}
\textbf{根因}：\gc{lmsysorg/sglang:latest}（2 天前构建）按 \textbf{CUDA 13.0} 编译，而本机驱动最高只支持 \textbf{CUDA 12.9}。
\vspace{1mm}

\textbf{解决}：换用 CUDA 12.9 镜像并在本地打成 \gc{latest}（compose 文件无需改）：
\begin{lstlisting}[language=]
docker pull lmsysorg/sglang:v0.5.16-cu129-runtime
docker tag  lmsysorg/sglang:v0.5.16-cu129-runtime lmsysorg/sglang:latest
docker image prune -f   # 清掉 cu130 悬空镜像，回收约 30GB
\end{lstlisting}
\centering\scriptsize
备注：0.5.10（论文版本）只剩 cu130 镜像，无 cu12x，故改用 v0.5.16-cu129 匹配 12.9 驱动。
\end{frame}

% ------------------------------------------------------------------
\begin{frame}[fragile]{问题 \cnum{2}：\gc{--enable-prefix-caching} 被新版移除}
\textbf{现象}：容器能启动，但立刻崩溃重启，日志报：
\begin{lstlisting}[language=]
sglang serve: error: unrecognized arguments: --enable-prefix-caching
\end{lstlisting}
\textbf{根因}：新版 SGLang v0.5.16 \textbf{移除了 \gc{--enable-prefix-caching}}——RadixAttention 前缀缓存现在\textbf{默认开启}，不再需要该 flag；而 compose 文件是照旧版 0.5.10 写的。
\vspace{1mm}

\textbf{解决}：从所有 compose 里删掉该参数（功能不受影响，前缀缓存仍默认开）：
\begin{lstlisting}[language=]
sed -i '/--enable-prefix-caching/d' \
  docker-compose.dp.yml docker-compose.tp.yml docker-compose.ep.yml \
  docker-compose.pd.yml docker-compose.pd-node.yml
\end{lstlisting}
\centering\scriptsize
两个坑都是\textbf{环境 / 版本问题}，与代码清理无关，且都已解决。
\end{frame}

% ------------------------------------------------------------------
\begin{frame}[fragile]{问题 \cnum{3}：多卡 NCCL 失败（TP=4 首次启动）}
\textbf{现象}：sglang-tp 容器反复重启（\gc{RestartCount=24}），日志：
\begin{lstlisting}[language=]
RuntimeError: NCCL error: unhandled system error   # ncclCommInitRank (4 卡建通信组)
\end{lstlisting}
\textbf{根因}：Docker 里多卡 NCCL 需要更大共享内存；且 RTX 4090 无 NVLink、PCIe 上 P2P 不稳。
（DP/PD 都是单卡容器 \gc{tp\_size=1}，不触发跨卡 NCCL，所以之前没事。）
\vspace{1mm}

\textbf{解决}：给 sglang-tp / sglang-ep 服务加 \gc{ipc: host} 与 NCCL 兜底环境变量：
\begin{lstlisting}[language=]
# docker-compose.tp.yml / ep.yml 的 sglang 服务加：
ipc: host
environment:
  - NCCL_P2P_DISABLE=1
  - NCCL_IB_DISABLE=1
\end{lstlisting}
\centering\scriptsize
代价：All-Reduce / All-to-All 走较慢的共享内存通路，延迟略高（这解释了 TP/EP 延迟比论文高）。
\end{frame}

% ------------------------------------------------------------------
\begin{frame}[fragile]{问题 \cnum{4}：\gc{--enable-expert-parallel} 被废弃（EP=4 启动）}
\textbf{现象}：sglang-ep 崩溃，日志报：
\begin{lstlisting}[language=]
sglang serve: error: unrecognized arguments: --enable-expert-parallel
\end{lstlisting}
\textbf{根因}：新版 SGLang v0.5.16 用 \gc{--ep-size N} 取代了 \gc{--enable-expert-parallel}。
\vspace{1mm}

\textbf{解决}：一条 \gc{sed} 替换即可：
\begin{lstlisting}[language=]
sed -i 's/--enable-expert-parallel/--ep-size 4/' docker-compose.ep.yml
\end{lstlisting}
\vfill
\begin{block}{网络与磁盘补充}
\scriptsize
\begin{itemize}
  \item \textbf{网络}：跑 TP/EP 期间本地到 VM 的 SSH 多次瞬断（VM 一直在运行，是网络路径问题）$\to$ 改用
  \textbf{VM 上 detached 脚本 + 短命令轮询}，大模型下载/加载不受 SSH 断影响。
  \item \textbf{磁盘}：每个大模型实验前先 \gc{docker volume rm radixgates\_hf\_cache} 腾空间——
  Qwen2.5-32B $\approx$ 64GB、MoE $\approx$ 28GB，119GB 盘装不下全部模型。
\end{itemize}
\end{block}
\end{frame}

% ==================================================================
\section*{六个实验}

\begin{frame}{问题汇总表（4 个坑，全是环境 / 版本问题）}
\scriptsize
\begin{tabular}{@{}p{1.9cm}p{3.0cm}p{3.4cm}p{3.2cm}@{}}
\toprule
\textbf{问题} & \textbf{现象} & \textbf{根因} & \textbf{解决} \\
\midrule
\cnum{1} 镜像 CUDA 版本 &
容器起不来，\gc{cuda>=13.0} 报错 &
\gc{sglang:latest} 按 CUDA 13.0 编译，驱动只支持 12.9 &
换 \gc{v0.5.16-cu129-runtime} 镜像并 retag 成 latest \\
\addlinespace
\cnum{2} 废弃参数 &
容器崩溃重启，\gc{unrecognized arguments} &
新版 SGLang 移除 \gc{--enable-prefix-caching}（前缀缓存默认开） &
用 \gc{sed} 从所有 compose 删除该参数 \\
\addlinespace
\cnum{3} 多卡 NCCL 失败 &
sglang-tp 反复重启，\gc{NCCL error: unhandled system error} &
Docker 多卡 NCCL 需更大共享内存；4090 无 NVLink、PCIe P2P 不稳 &
加 \gc{ipc: host} + \gc{NCCL\_P2P\_DISABLE=1} / \gc{NCCL\_IB\_DISABLE=1} \\
\addlinespace
\cnum{4} EP 参数被废弃 &
sglang-ep 崩溃，\gc{unrecognized arguments: --enable-expert-parallel} &
v0.5.16 用 \gc{--ep-size N} 取代旧 flag &
\gc{sed} 替换成 \gc{--ep-size 4} \\
\bottomrule
\end{tabular}
\vfill
\begin{itemize}\scriptsize
  \item 另一个非阻塞小瑕疵：demo 的 \gc{\_log\_utils} 硬件检测把本机 PCIe 误标成 "NVLink"，
  只影响日志文件名标签，\textbf{不影响结果与数值}。
\end{itemize}
\end{frame}

% ------------------------------------------------------------------
\begin{frame}{六个实验总览：目的与对应生产场景}
\scriptsize
\begin{tabular}{@{}p{2.4cm}p{4.9cm}p{4.9cm}@{}}
\toprule
\textbf{实验} & \textbf{验证目的} & \textbf{对应生产场景} \\
\midrule
\cnum{1} 流式 SSE &
token 逐字实时流回，验证 direct SSE 链路 &
实时对话 / 代码助手的流式输出 \\
\addlinespace
\cnum{2} PD 分离 &
Prefill/Decode 分卡 + mooncake 传 KV，30 并发低 TTFT &
金融实时风控 / 低延迟推理 \\
\addlinespace
\cnum{3} DP=4 &
4 副本 + 前缀哈希路由，40 并发、4$\times$ 扩展、KV 命中 &
企业内网 AI 助手（数百员工并发、按部门路由） \\
\addlinespace
\cnum{4} 企业问答压测 &
多部门并发 + 信号量背压，验证生产稳定性 &
企业问答系统生产验收 \\
\addlinespace
\cnum{5} TP=4 &
32B 权重切 4 卡，每层 All-Reduce &
金融合规大模型（32B+） \\
\addlinespace
\cnum{6} EP=4 &
MoE 的 Expert 分卡 + All-to-All &
医疗 / 法律 MoE 模型（DeepSeek-V3 类） \\
\bottomrule
\end{tabular}
\vfill
\centering\scriptsize 实验一$\sim$三 用 Llama-3-8B；实验五 用 Qwen2.5-32B；实验六 用 Qwen1.5-MoE-A2.7B。
\end{frame}

% ------------------------------------------------------------------
\begin{frame}[fragile]{实验一 \cnum{1}：单卡 / 流式 SSE（streaming\_demo）}
\textbf{目的}：演示 token 逐字实时流回（打字机效果），验证 Gateway$\to$SGLang 的 direct SSE 链路。\\
\textbf{对应生产}：实时对话 / 代码助手的流式输出。
\vspace{1mm}
\begin{lstlisting}[language=bash]
# 脚本化（非交互）：把问题用管道喂进去
echo "用一句话解释什么是GPU" | python3 examples/streaming_demo.py
\end{lstlisting}
\begin{block}{关键结果}
\small
\begin{itemize}
  \item Gateway：\gc{direct} 模式（无 Kafka、无 WebSocket，纯 SSE）
  \item 输出 \textbf{$\sim$58 tokens}，耗时 \textbf{1.04s}，吞吐 \textbf{$\sim$56.0 tokens/s}
\end{itemize}
\end{block}
\end{frame}

% ------------------------------------------------------------------
\begin{frame}[fragile]{实验二 \cnum{2}：PD 分离（pd\_demo）}
\textbf{目的}：Prefill(GPU0) / Decode(GPU1) 分卡 + mooncake 传 KV，验证 30 并发与低 TTFT。\\
\textbf{对应生产}：金融实时风控 / 低延迟推理。
\vspace{1mm}
\begin{lstlisting}[language=bash]
docker compose -f docker-compose.dp.yml down     # 先释放 4 卡
docker compose -f docker-compose.pd.yml up -d    # prefill+decode+router+gateway
curl -s localhost:8080/health                    # {"status":"ok","mode":"direct"}
python3 examples/pd_demo.py
\end{lstlisting}
\begin{block}{关键结果}
\small
\begin{itemize}
  \item 并发吞吐（30 并发）：\textbf{成功 30/30}，总耗时 \textbf{5.1s}，\textbf{P50=2544ms ｜ P95=4973ms}
  \item TTFT：short 94.7ms ｜ medium 93.9ms ｜ long 117.3ms（warm TTFT $\sim$90ms）
  \item disaggregation 参数（mooncake / router \gc{--pd-disaggregation}）在 v0.5.16 完全兼容
\end{itemize}
\end{block}
\end{frame}

% ------------------------------------------------------------------
\begin{frame}[fragile]{实验三 \cnum{3}：DP=4 数据并行（dp\_demo）}
\textbf{目的}：4 个独立副本 + 前缀哈希一致性路由，验证 40 并发、4$\times$ 扩展、KV Cache 命中。\\
\textbf{对应生产}：企业内网 AI 助手（数百员工并发、按部门路由）。
\vspace{1mm}
\begin{lstlisting}[language=bash]
MODEL_PATH=NousResearch/Meta-Llama-3-8B-Instruct \
  docker compose -f docker-compose.dp.yml up -d   # 4 副本 + 网关
python3 examples/dp_demo.py
\end{lstlisting}
\begin{block}{关键结果}
\small
\begin{itemize}
  \item 4/4 节点健康；四卡对称显存 \textbf{21246/24564 MB (86.5\%)}，无权重切分
  \item 并发吞吐（40 并发）：\textbf{成功 40/40}，\textbf{P50=1215ms ｜ P95=2300ms}，加速比 \textbf{$\sim$4.0$\times$}
  \item 冷 / 热 TTFT（prefix-0）：\textbf{warm 41.5ms vs cold 89.6ms} $\Rightarrow$ \textbf{Cache Hit}
\end{itemize}
\end{block}
\end{frame}

% ------------------------------------------------------------------
\begin{frame}[fragile]{实验四 \cnum{4}：企业问答并发压测（enterprise\_qa\_demo）}
\textbf{目的}：多部门并发 + 信号量背压（生产稳定性）。\\
\textbf{对应生产}：企业问答系统生产验收（跑在当前 PD 网关上，与后端模式无关）。
\vspace{1mm}
\begin{lstlisting}[language=bash]
python3 examples/enterprise_qa_demo.py
\end{lstlisting}
\begin{block}{关键结果}
\small
\begin{itemize}
  \item 30 名员工、5 个部门同时提问
  \item \textbf{成功 30/30}，总耗时 \textbf{10.2s}，\textbf{零失败}
  \item “信号量限流完美运作：所有请求在等待池中有序处理，零丢失”
\end{itemize}
\end{block}
\end{frame}

% ------------------------------------------------------------------
\begin{frame}[fragile]{实验五 \cnum{5}：TP=4 张量并行（tp\_demo，Qwen2.5-32B）}
\textbf{目的}：单卡放不下的 32B 权重切 4 卡，每层 All-Reduce。\\
\textbf{对应生产}：金融合规大模型（32B+ 单卡放不下）。
\vspace{1mm}
\begin{lstlisting}[language=bash]
CUDA_VISIBLE_DEVICES=0,1,2,3 TP_SIZE=4 \
  MODEL_PATH=Qwen/Qwen2.5-32B-Instruct \
  docker compose -f docker-compose.tp.yml up -d
python3 examples/tp_demo.py
\end{lstlisting}
\begin{block}{关键结果}
\small
\begin{itemize}
  \item 日志确认 \gc{tp\_size=4}；四卡对称显存 \textbf{22320/24564 MB (90.9\%)} $\Rightarrow$ 32B 已切 4 卡
  \item 并发（20）：\textbf{成功 20/20}，\textbf{P50=4836ms ｜ P95=6475ms}
  \item 对照 README（3851ms）偏高 $\sim$25\%：4090 无 NVLink，\gc{NCCL\_P2P\_DISABLE} 走较慢 PCIe All-Reduce；结构完全复现
\end{itemize}
\end{block}
\end{frame}

% ------------------------------------------------------------------
\begin{frame}[fragile]{实验六 \cnum{6}：EP=4 专家并行（ep\_demo，Qwen1.5-MoE）}
\textbf{目的}：MoE 的 Expert 分卡 + token 路由 All-to-All。\\
\textbf{对应生产}：医疗 / 法律 MoE 模型（DeepSeek-V3 类）。
\vspace{1mm}
\begin{lstlisting}[language=bash]
CUDA_VISIBLE_DEVICES=0,1,2,3 TP_SIZE=4 \
  MODEL_PATH=Qwen/Qwen1.5-MoE-A2.7B-Chat \
  docker compose -f docker-compose.ep.yml up -d
python3 examples/ep_demo.py
\end{lstlisting}
\begin{block}{关键结果}
\small
\begin{itemize}
  \item 日志确认 \gc{ep\_size=4} + \textbf{4 个 EP rank}（\gc{[TP0 EP0]}$\sim$\gc{[TP3 EP3]}）；显存 91.0\%，Expert 分片
  \item 并发（20）：\textbf{成功 20/20}，\textbf{P50=1653ms ｜ P95=2095ms}
  \item 对照 README（1099ms）偏高：同因 PCIe All-to-All 较慢；结构完全复现
\end{itemize}
\end{block}
\end{frame}

% ==================================================================
\section*{对照与结论}

\begin{frame}{结果对照表：六实验 vs README 基准（4090 列）}
\scriptsize
\begin{tabular}{@{}p{2.0cm}p{4.6cm}p{3.4cm}p{1.5cm}@{}}
\toprule
\textbf{实验} & \textbf{本次实测} & \textbf{README 基准} & \textbf{吻合} \\
\midrule
\cnum{1} 流式 SSE & 58 tokens，$\sim$56 tok/s，1.04s & —（功能演示） & \textcolor{rggreen}{$\checkmark$} \\
\addlinespace
\cnum{2} PD 分离 & 30/30，5.1s，P50 2544ms，P95 4973ms & 30/30，5.2s，P50 2.8s，P95 5.2s & \textcolor{rggreen}{$\checkmark$} 吻合 \\
\addlinespace
\cnum{3} DP=4 & 40/40，P50 1215ms，$\sim$4.0$\times$ & 40/40，P50 1174ms，4$\times$ & \textcolor{rggreen}{$\checkmark$} 吻合 \\
\addlinespace
\cnum{4} 企业问答 & 30/30，10.2s，信号量零丢失 & —（稳定性验收） & \textcolor{rggreen}{$\checkmark$} \\
\addlinespace
\cnum{5} TP=4 & 20/20，P50 4836ms，tp\_size=4，32B 切 4 卡 & 20/20，P50 3851ms & \textcolor{rgorange}{$\checkmark$}* 结构 \\
\addlinespace
\cnum{6} EP=4 & 20/20，P50 1653ms，ep\_size=4，4 EP rank & 20/20，P50 1099ms & \textcolor{rgorange}{$\checkmark$}* 结构 \\
\bottomrule
\end{tabular}
\vfill
\centering\scriptsize
*TP/EP 延迟偏高约 25–50\%：4090 无 NVLink、\gc{NCCL\_P2P\_DISABLE} 走较慢 PCIe 通信；成功率 / 分片 / EP rank 全部正确。
\end{frame}

% ------------------------------------------------------------------
\begin{frame}{结论要点}
\begin{itemize}\small
  \item \textbf{六个实验全部跑通}；DP / PD / TP / EP 四种并行的\textbf{结构性结果全部复现 README}。
  \item TP / EP 延迟比论文高约 25–50\%：这台 4090 无 NVLink，需 \gc{NCCL\_P2P\_DISABLE} 兜底走较慢 PCIe 通信；但\textbf{成功率、权重分片、EP rank 全部正确}。
  \item 这是\textbf{清理后的纯 direct 架构}（已删除全部 Kafka / Redis / C++ 旧代码）——实验证明清理\textbf{没有破坏任何功能}：\gc{\{"status":"ok","mode":"direct"\}}，前缀路由 + 信号量 + SSE + PD 分离全部正常。
  \item 全程踩了 \textbf{4 个坑}（\cnum{1} 镜像 CUDA 版本、\cnum{2} \gc{--enable-prefix-caching}、\cnum{3} 多卡 NCCL、\cnum{4} \gc{--ep-size}），\textbf{全是环境 / 版本问题，与代码清理无关}，且都已解决。
\end{itemize}
\vfill
\begin{block}{一句话总结}
\small 在 4$\times$ RTX 4090（PCIe）上，用 v0.5.16-cu129 复现了 RadixGates 的 DP / PD / TP / EP 四种并行 + 流式 + 企业问答，
结果与论文基准一致，纯 direct 架构功能完好。
\end{block}
\vfill
\centering {\large\textbf{谢谢！}}
\end{frame}

\end{document}
